\section{Evaluation} \label{sec:evaluation}

% the number of copyright holder

In this section, we present the evaluation results of \tool on real-world datasets and answer the following research questions. % Specifically, we present the experimental results of \tool in three phases, i.e., license understanding, incompatibility issue localization, and license incompatibility resolution, and answer the following research 

    \textbf{RQ1}: Can \tool provide a fine-grained understanding of license texts by effectively extracting entities and their relations?
    
    \textbf{RQ2}: Can \tool outperform state-of-the-art approaches in detecting and localizing license incompatibility issues?
    
    \textbf{RQ3}: How effective is \tool for resolving license incompatibility issues in real-world OSS?


\iffalse 

\subsection{Data Preparation}


\subsubsection{\textbf{Dataset of License Entities and Relations.}}

The license understanding phase includes two tasks, i.e., entity extraction and relation extraction. For both tasks we performed labelling, training, and testing in \textbf{sentences}. The collection of the sentence datasets includes three steps: (1) \textbf{License collection}. We collected license texts from TLDRLegal~\cite{tldrlegal}, which provides a quick summary, full-text, and license terms associated with attitudes for 212 licenses. \sihan{In addition, we also utilized the datasets provide by \cite{xu2021lidetector} which  XXX projects on GitHub and obtained \sihan{XXX} licenses including official and custom ones. }(2) \textbf{Sentence split.} For the collected licenses, we split them into sentences with CoreNLP~\cite{url-corenlp}, and then filtered out duplicated licenses and sentences. In total, we obtained 900 unique sentences for training, validation, and testing. (3) \textbf{Data labelling}. Given a license sentence, we labeled the entities and relations in the training phase. Specifically, we labeled four types of entities, i.e., \textit{action}, \textit{object}, \textit{attitude}, and \textit{condition} to train the entity extraction model. We also labeled five types of the relations between these entities, i.e., \textit{action-object}, \textit{action-attitude}, \textit{action-condition}, \textit{condition-action}, and \textit{other} to train the relation extraction model. Note that \textit{action-condition} denotes an \textit{action} under a \textit{condition}, while \textit{condition-action} denotes the \textit{action} in the \textit{condition} clause. 
For \textit{actions} and \textit{attitudes}, we verified their labels by comparing them with the datasets in previous studies~\cite{tldrlegal,xu2021lidetector}. For \textit{objects} and \textit{conditions}, the labels were cross-validated by three authors. 
% We randomly split these samples into the training and testing datasets by 4:1, and further split the training dataset into the training and validation datasets by 4:1. 
{We randomly split these samples into the training, validation and testing datasets by 3:1:1.} 
The statistics of these datasets can be seen in \Cref{tab:dataset}.

\begin{table}
	\centering
	\small
	\caption{Statistics of Datasets for License Understanding \sihan{to update}}
	\scalebox{0.9}{\begin{tabular}{cc|cc}
		\toprule
		%\hline
		\textbf{Entity} & \textbf{No.} & \textbf{Relation} & \textbf{No.}\\
		\hline
		action  & XX (xx\%) & action-object & XX (xx\%) \\
		object  & XX (xx\%) & action-attitude & XX (xx\%) \\
		attitude  & XX (xx\%) & action-condition & XX (xx\%) \\
		condition  & XX (xx\%) & condition-action & XX (xx\%) \\
		\hline
		Total  & XX (xx\%) & Total & XX (xx\%) \\
		%\hline
		\bottomrule
		\label{tab:dataset}
	\end{tabular}}
\end{table}







\subsubsection{\textbf{Dataset of OSS Projects.}}

To evaluate the performance of \tool on localizing and resolving license incompatibility issues, we crawled \sihan{XXX} popular OSS projects with more than \sihan{1,000} stars on GitHub, from which we randomly selected \sihan{XXX} projects to evaluate the performance of \tool in localizing and resolving license incompatibility issues. To construct the ground-truth datasets, three authors cross-validated the labels and a lawyer was involved in the validation. We have made the dataset publicly available~\cite{ourwebsite}.


\fi




\subsection{Evaluation on License Understanding}
In this section, we evaluate the performance of \tool on fine-grained entity extraction and relation extraction.
%\subsubsection{\textbf{Setup}}  
For both tasks we performed labelling, training, and testing in \textbf{sentences}. The sentence datasets were constructed based on the dataset from a previous study~\cite{xu2021lidetector}. \revise{To reduce bias, we did not set aside any licenses when evaluating the effectiveness of license understanding provided by \tool. In total, there were 212 official licenses and 188 custom licenses, comprising 21,844 license sentences  (i.e., {11,973} from \textit{tldrlegal}~\cite{tldrlegal} and {9,871} from GitHub). We manually labelled the entities and relations. During labeling, we took Fleiss' Kappa to assess the reliability of agreement among the raters. The result is 0.83, indicating that the raters reach high agreement.} We randomly split each dataset into the training, validation and testing datasets by 3:1:1. 

\subsubsection{{Entity Extraction}.}
% To evaluate the performance of \tool for entity extraction, we compare it with two state-of-the-art tools (i.e., FOSS-LTE~\cite{APSEC`17-Kapitsaki-termsIdentifying} and LiDetector~\cite{xu2021lidetector}) and two natural language processing techniques (i.e., regular matching~\cite{USENIX`19-Andow-Policylint} and semantic similarity~\cite{goldberg2014word2vec}) as follows:

% \begin{itemize}[leftmargin=*]
%     \item \textbf{FOSS-LTE}~\cite{APSEC`17-Kapitsaki-termsIdentifying}, which employed a topic model named Latent Dirichlet Allocation (LDA)~\cite{MTA`19-Jelodar-LDA} to identify license terms.
%     \item \textbf{LiDetector}~\cite{xu2021lidetector}, which utilized sequence labelling and sentiment analysis to identify license term entities (similar to the \textit{actions} in this paper) and \textit{attitudes}, respectively.
%     \item \textbf{Regular matching}~\cite{USENIX`19-Andow-Policylint}, which was implemented by predefining a set of keyword patterns for entity extraction. 
%     \item \textbf{Semantic similarity}~\cite{goldberg2014word2vec}, which utilized the word2vec pre-trained language model~\cite{mikolov2013distributed,mikolov2013efficient} to measure the cosine similarity between the tokens and predefined keywords. 
% \end{itemize}

To evaluate the performance of \tool for entity extraction, we compare it with two state-of-the-art tools (i.e., \textbf{FOSS-LTE}~\cite{APSEC`17-Kapitsaki-termsIdentifying} and \textbf{LiDetector}~\cite{xu2021lidetector}) and two natural language processing techniques (i.e., \textbf{regular matching}~\cite{USENIX`19-Andow-Policylint} and \textbf{semantic similarity}~\cite{goldberg2014word2vec}).
\revise{Specifically, {FOSS-LTE}~\cite{APSEC`17-Kapitsaki-termsIdentifying} employed a topic model named Latent Dirichlet Allocation (LDA)~\cite{MTA`19-Jelodar-LDA} to identify license terms.
{LiDetector}~\cite{xu2021lidetector} utilized sequence labelling and sentiment analysis to identify license term entities (similar to the \textit{actions} in this paper) and {attitudes}.
{Regular matching}~\cite{USENIX`19-Andow-Policylint} was implemented by pre-defining a set of keyword patterns for entity extraction.
{Semantic similarity}~\cite{goldberg2014word2vec} utilized the word2vec pre-trained language model~\cite{mikolov2013distributed,mikolov2013efficient} to measure the cosine similarity between the tokens and predefined keywords.}
%
Note that the outputs of FOSS-LTE are license term phrases mapped from topic sentences (e.g., MayGrantPatents). To conduct a fair comparison of license understanding, we split the outputs of FOSS-LTE into four types of entities as other baselines. Similarly, LiDetector naturally extracts license terms (e.g., Commercial Use), attitudes (e.g., CANNOT), and conditions (e.g., \textit{provided that}) from license texts, which can be compared with the \textit{actions}, \textit{attitudes}, and \textit{conditions} in this paper. As for regular matching and semantic similarity, we analyzed 900 license sentences and summarized a set of keyword patterns for each type of entities as shown online~\cite{ourwebsite}. 

%\subsubsection{\textbf{Results}}

\begin{table}[]
\centering
\small
\caption{Comparison on Fine-grained Entity Extraction}
\scalebox{0.73}{\begin{tabular}{ccccccccccc}
\toprule
\multirow{2}{*}{\textbf{Tool}} & \multicolumn{2}{c}{\textbf{Actions}}              & \multicolumn{2}{c}{\textbf{Objects}}           & \multicolumn{2}{c}{\textbf{Attitudes}}             & \multicolumn{2}{c}{\textbf{Conditions}}    & \multicolumn{2}{c}{{\revise{\textbf{Total}}}}       \\\cline{2-11}
                                 & \textbf{P (\%)} & \textbf{R (\%)} & \textbf{P (\%)} & \textbf{R (\%)} &  \textbf{P (\%)} & \textbf{R (\%)} &  \textbf{P (\%)} & \textbf{R (\%)} & \revise{\textbf{P (\%)}} & \revise{\textbf{R (\%)}} \\
\midrule
{Reg. Match.} & 54.76             & 65.07             & 5.32            & 22.38             & 2.58             & 47.61             & 46.87             & 61.64      & \revise{14.74} & \revise{49.41}        \\
Similarity              & 50.81             & 59.96            & 3.18             & 12.27             & 2.35             & 40.47             & 46.86             & 61.65      & \revise{14.28} & \revise{42.33}       \\
FOSS-LTE                         & 35.98             & 23.12             & 5.76             & 19.13             & 3.06             & 26.19             & 66.67            & 21.92       & \revise{14.30} & \revise{21.12}      \\
LiDetector                       & 43.14             & 55.79             & -             & -            & 74.74             & 56.35             & {90.25}             & 86.30         & \revise{47.79} & \revise{32.80}     \\
\textbf{\tool}                             & \textbf{80.79}             & \textbf{84.86}             & \textbf{59.23}             & \textbf{63.83}             & \textbf{78.69}             & \textbf{76.80}             & \textbf{95.45}             & \textbf{87.50}         & \textbf{\revise{74.15}} & \textbf{\revise{77.31}}     \\
\bottomrule
\label{tab:entity}
\end{tabular}}
\end{table}

\Cref{tab:entity} displays the experimental results of entity extraction for five approaches. Note that since LiDetector does not extract the \textit{objects} of \textit{actions} in license texts, the results of extracting \textit{objects} by LiDetector are denoted by ``-''. In total, \tool achieves 74.07\% precision and 77.36\% recall for entity extraction. It can be seen that the precision of \tool is much higher than that of the baselines by at least 26.03\%, 53.47\%, 3.95\%, and 5.2\% when extracting \textit{actions}, \textit{objects}, \textit{attitudes}, and \textit{conditions}, respectively. It indicates that \tool can precisely extract entities in license texts without much noise. However, regular matching and similarity-based approaches achieve low precision especially for \textit{objects} and \textit{attitudes}. By analyzing the results, a possible explanation could be the ambiguous expressions for the two types of entities. 
Moreover, the recall scores achieved by \tool are also higher than those of four baselines by at least 19.79\%, 41.45\%, 29.19\%, and 1.2\% when extracting \textit{actions}, \textit{objects}, \textit{attitudes}, and \textit{conditions}, respectively. %\tool achieved the highest performance on \textit{conditions}. A possible reason is that the expressions of \textit{conditions} are less flexible than other types of entities. 
Among four baselines, LiDetector achieves a competitive performance with \tool on extracting \textit{attitudes} and \textit{conditions}. 
However, the precision and recall scores of LiDetector on extracting \textit{actions} are lower than those of \tool. By analyzing the entities extracted by LiDetector, we found that {the entities extracted by LiDetector are more coarse-grained, sometimes along with its \textit{objects} or \textit{attitudes}, which made them longer than the entities extracted by other approaches, which decreased the evaluation results of LiDetector.}

\iffalse
\begin{table*}[]
\centering
\small
\caption{Comparison on Entity Extraction\sihan{ to update}}
\scalebox{0.9}{\begin{tabular}{cccccccccccccccc}
\toprule
\multirow{2}{*}{\textbf{Tool}} & \multicolumn{3}{c}{\textbf{Actions}}              & \multicolumn{3}{c}{\textbf{Objects}}           & \multicolumn{3}{c}{\textbf{Attitudes}}             & \multicolumn{3}{c}{\textbf{Conditions}}           & \multicolumn{3}{c}{\textbf{Total}}                \\\cline{2-16}
                                 & \textbf{P (\%)} & \textbf{R (\%)} & \textbf{F1 (\%)} & \textbf{P (\%)} & \textbf{R (\%)} & \textbf{F1 (\%)} & \textbf{P (\%)} & \textbf{R (\%)} & \textbf{F1 (\%)} & \textbf{P (\%)} & \textbf{R (\%)} & \textbf{F1 (\%)} & \textbf{P (\%)} & \textbf{R (\%)} & \textbf{F1 (\%)} \\
\midrule
Regular Matching                 & 50.43             & 74.32             & 60.09              & 5.20            & 25.27             & 8.62              & 2.15             & 42.06             & 4.09             & 54.55             & 98.63             & 70.24              & 14.26             & 55.58             & 22.71              \\
Semantic Similarity              & 47.36             & 67.15            & 55.54              & 3.19             & 11.91             & 5.03              & 1.98             & 36.51             & 3.76             & 54.55             & 98.63             & 70.24              & 13.57             & 44.33             & 20.78              \\
FOSS-LTE                         & 35.98             & 23.12             & 28.15              & 5.76             & 19.13             & 8.86              & 3.06             & 26.19             & 5.48              & 66.67            & 21.92             & 32.98              & 14.30             & 21.12             & 17.06              \\
LiDetector                       & XX             & XX             & XX              & -             & -             & -              & XX             & XX             & XX              & XX             & XX             & XX              & XX             & XX             & XX              \\
\textbf{\tool}                             & \textbf{70.00}             & \textbf{72.80}             & \textbf{71.37}              & \textbf{42.07}             & \textbf{47.27}             & \textbf{44.52}              & \textbf{76.61}             & \textbf{75.40}             & \textbf{76.00}              & \textbf{95.83}             & \textbf{94.52}             & \textbf{95.17}              & \textbf{64.85}             & \textbf{68.15}             & \textbf{66.46}            \\
\bottomrule
\label{tab:entity}
\end{tabular}}
\end{table*}
\fi






%\noindent $\bullet$ \textbf{Relation extraction}. 
\subsubsection{{Relation Extraction}}
\Cref{fig:relation} shows the performance of two approaches on relation extraction. Since FOSS-LTE~\cite{APSEC`17-Kapitsaki-termsIdentifying} and LiDetector~\cite{xu2021lidetector} do not extract relations between entities, we did not compare \tool with them. Instead, we implemented a heuristic approach based on the positions and types of entities (denoted by \textbf{PTE}). Specifically, it determined entity relations according to their types and relative positions in the same sentence. For instance, \textit{action-object} indicates an \textit{action} in front of \textit{object}. Similarly, \textit{action-condition} and \textit{condition-action} were also distinguished by the relative positions of the \textit{action} and \textit{condition}. %As for grammar parsing, it first utilized CoreNLP~\cite{url-corenlp} for dependency parsing, which provided semantic dependencies between phrases in a sentence. Then, it classified the relations between entities according to these dependencies. For example, the relation between a pair of entities with the \textit{obj} dependency is considered as \textit{action-object}. The detailed mappings between semantic dependencies and entity relations are available online~\cite{ourwebsite}.
Finally, to conduct a fair comparison, we fed both approaches the same pairs of entities as input. 
%
%In \Cref{tab:realtion}, we use \textit{Act}, \textit{Obj}, \textit{Att}, and \textit{Con} to represent \textit{action}, \textit{object}, \textit{attitude}, and \textit{condition}, respectively. 
Since the input of entity relation is a pair of entities, there exist no false negatives and thus we evaluate the performance of two approaches only with the accuracy metric. It can be seen that \tool can effectively identify entity relations, achieving 95.69\% accuracy in total. We can also observe that the performance of \tool is superior to that of the heuristic baseline for all types of relations. The baseline only achieves comparable performance with \tool on \textit{action-condition} and \textit{condition-action}. It indicates that the relation between an \textit{action} and a \textit{condition} can be easier inferred by their relative positions in the sentence compared to other types of entity relations. However, since PTE requires the types of entities as input, its performance relies on the accuracy of entity extraction, while \tool has no requirements for such information.

%In detail, a pair of entities with \textit{obj} dependency is considered as \textit{action-recipient}, a pair of entities with \textit{aux, advmod} dependencies is considered as \textit{action-attitude}, and a pair of entities with \textit{acl, advcl} dependencies (usually from \textit{action} to \textit{condition}) is considered as existing \textit{condition} relation, the dependency direction of which decides it falls into \textit{action-condition} or \textit{condition-action}. 

\noindent \fbox{
\parbox{0.95\linewidth}{
\textbf{Answer to RQ1:} 
\tool can effectively extract fine-grained entities and their relations from license texts, with 74.07\% precision and 77.36\% recall for entity extraction, 95.09\% accuracy for relation extraction, which provides detailed information for license incompatibility localization and resolution.
}}


\iffalse
\begin{table}
	\centering
	\small
	\caption{Accuracy of Relation Extraction \sihan{to update}}
	%\scalebox{0.9}{
    \begin{tabular}{ccccc|c}
		\toprule
		%\hline
		\textbf{Method} & \textbf{Act-Obj} & \textbf{Act-Att} & \textbf{Act-Con} & \textbf{Con-Act} & \textbf{Total}\\
		\hline
		 PTE  & 87.43  & 69.09 & 87.51 & 99.22 & 74.06\\
		%Grammar Parsing  & 42.77  & 29.29 & 1.01 & 0.00 & 16.80\\
		\textbf{\tool}  & \textbf{98.46}  & \textbf{93.87} & \textbf{86.88} & \textbf{100.00} & \textbf{95.69}\\
		%\hline
		\bottomrule
		\label{tab:realtion}
	\end{tabular}%}
\end{table}
\fi

\begin{figure}[]
    \centering
    \includegraphics[width=0.98\linewidth]{figs/fig-relation-5.pdf}
    %\vspace{-2mm}
    \caption{Accuracy of Relation Extraction}
    \label{fig:relation}
\end{figure}



\iffalse
\begin{table}
	\centering
	\small
	\caption{Comparison on License Understanding \sihan{to update}}
	\scalebox{0.9}{\begin{tabular}{ccccc}
		\toprule
		%\hline
		\textbf{Entity Extraction} & \textbf{Relation Extraction}& P & R & C\\
		\midrule
		Regular Matching  & Regular Matching  &XX  & XX & XX \\
		Semantic Similarity  & Regular Matching  & XX  & XX & XX \\
		FOSS-LTE  & FOSS-LTE  & XX  & XX & XX \\
		LiDetector  & LiDetector  & XX  & XX & XX \\
		\textbf{\tool}  &\textbf{\tool}  & \textbf{XX}  & \textbf{XX} & \textbf{XX} \\
		%\hline
		\bottomrule
		\label{tab:realtion}
	\end{tabular}}
\end{table}
\fi

\subsection{Evaluation on Incompatibility Localization}\label{sec:incomp}
%All the techniques and tools were evaluated on the same testing dataset. 
%In this section, we evaluate the ability of \tool on detecting and localizing license incompatibility issues.
\subsubsection{{Setup}}
We evaluate the performance of \tool on 300 real-world projects on GitHub. Specifically, we first crawled 10,000 popular projects with more than 1,000 stars. Then, we randomly selected 300 projects and constructed a ground-truth dataset. Three authors and a lawyer were involved in labelling and verifying the dataset which has been made publicly available online~\cite{ourwebsite}. In total, we extracted {81,402} licenses for incompatibility analysis. %, including \sihan{XXX} unique official licenses and \sihan{XXX} custom licenses.  
We compare \tool with four state-of-the-art tools, i.e., the SPDX Violation Tools (\textbf{SPDX-VT})~\cite{jss`17-Kapitsaki-SPDX}, \textbf{Ninka}~\cite{ase`10-German-Ninka} equipped with \textit{tldrlegal}~\cite{tldrlegal}, {\textbf{Librariesio}~\cite{url-librariesio}, and \textbf{LiDetector}~\cite{xu2021lidetector}}. {SPDX-VT}~\cite{jss`17-Kapitsaki-SPDX}, which designed a directed graph to represent the compatibility relations between a set of licenses in Software Package Data Exchange (SPDX)~\cite{spdx}. {Ninka}~\cite{ase`10-German-Ninka}, which identified licenses with a sentence-matching method. Combined with the summarized rights and obligations provided by \textit{tldrlegal}~\cite{tldrlegal}, Ninka can be utilized to detect incompatibilities between licenses. {Librariesio}~\cite{url-librariesio}, which checked compatibility between a predefined set of SPDX licenses. {LiDetector}~\cite{xu2021lidetector}, which firstly introduced a learning-based method to automatically identify license terms from arbitrary licenses and detect license incompatibilities. 
\iffalse
\begin{itemize}[leftmargin=*]
    \item \textbf{SPDX-VT}~\cite{jss`17-Kapitsaki-SPDX}, which designed a directed graph to represent the compatibility relations between a set of licenses in Software Package Data Exchange (SPDX)~\cite{spdx}.
    \item \textbf{Ninka}~\cite{ase`10-German-Ninka}, which identified licenses with a sentence-matching method. Combined with the summarized rights and obligations provided by \textit{tldrlegal}~\cite{tldrlegal}, Ninka can be utilized to detect incompatibilities between licenses.
    \item \textbf{Librariesio}~\cite{url-librariesio}, which checked compatibility between a predefined set of SPDX licenses. 
    \item \textbf{LiDetector}~\cite{xu2021lidetector}, which firstly introduced a learning-based method to automatically identify license terms from arbitrary licenses and detect license incompatibilities. 
\end{itemize}
\fi



\subsubsection{{Results}}



\Cref{tab:incom} displays the experimental results of license incompatibility localization. The results show the number of incompatible projects and incompatible issues reported by each approach. Note that in this paper, \textbf{an incompatible issue} represents a pair of licenses which are incompatible with each other and should be resolved to reduce legal risks. We also report the false positive (FP) rate and the false negative (FN) rate from both the project-level and the issue-level for five approaches. 
From~\Cref{tab:incom}, it can be seen that \tool is superior to all the baselines in terms of both FP and FN rates. Specifically, \tool reported 231 incompatible projects with 0.43\% FP rate  and no false negatives, which demonstrated its ability to detect projects with license incompatibility. 
% \gy{There are also other tools with the project-level FP to be zeros, i.e., Ninka equipped with \textit{tldrlegal} and SPDX-VT. It indicates that the projects reported by them are all incompatible projects without any false positives. However, their false negatives are also much higher than other approaches. The FN rates of Ninka and SPDX-VT are 42.86\% and 44.54\%, respectively. }

It can also be seen that Ninka equipped with \textit{tldrlegal} and SPDX-VT have low FP rates (i.e., 2.94\% and 2.63\%, respectively) but high FN rates (i.e., 42.61\% and 67.83\%).
By analyzing their detection processes and results, we found that both of them complied with strict and predefined rules for license compatibility analysis, which contributed to the low FP rates. However, relying on license identification and manually defined compatibility rules, the applications of these two approaches were limited to a small set of official licenses, while ignoring other licenses in the projects, which causes their high FN rates.
Similarly, the performance of Librariesio is close to that of Ninka and SPDX-VT.
% , with less false negatives but more false positives. 
It categorized {35} official licenses into five classes (e.g., \textit{permissive licenses}), and manually defined a heuristic rule based on the classification. It achieved a lower FN rate (i.e., 33.91\% at the project-level) than Ninka and SPDX-VT, since it took more licenses into consideration. However, due to the course-fined classification of licenses, it achieved the highest project-level FP rate among all the approaches (i.e., 11.11\%). 

Among four baselines, LiDetector achieved the closest performance to \tool for resolving license incompatibilities. Specifically, it reported 228 incompatible projects with 8 false positives and 10 false negatives, 
% By analyzing the false positive, we found that it occurred because \sihan{XXX}.
mainly caused by its coarse-grained license term extraction which did not distinguish between the same \textit{actions} with different \textit{objects} (e.g., \textit{Distribute source code} \textit{vs.} \textit{Distribute binaries}). Moreover, like other baselines, it also ignored the license hierarchy of the projects. 

\begin{table}[]
	\centering
	\footnotesize
	\caption{Comparison on License Incompatibility Localization} \label{tab:incom}
	\scalebox{0.95}{
	\begin{tabular}{cccc|ccc}
		\toprule
		%\hline
		%\textbf{Tool} & \textbf{\#Pro.} & \textbf{\#Incomp. Pro.} & \textbf{\#Incomp. Issue} & \textbf{FP} & \textbf{FN}\\
		\multirow{2}{*}{\textbf{Tool}} & \multicolumn{3}{c|}{\textbf{\#Incomp. Projects}} & \multicolumn{3}{c}{\textbf{\#Incomp. Issues}}\\\cline{2-7}
		 & {\textbf{\#Report.}} & \textbf{FP (\%)} & \textbf{FN (\%)} & \textbf{\#Report.}& \textbf{FP (\%)} & \textbf{FN (\%)}\\
		\hline
		Ninka   & 136 & 2.94 & 42.61 & 32,289 & 86.46 & 86.31 \\
		SPDX-VT   & 76 & 2.63 & 67.83 & 54,530 & 65.33 & 99.28 \\
		Librariesio   & 171 & 11.11 & 33.91 & 18,874 & 36.22 & 68.32 \\
		LiDetector   & 228 & 3.51 & 4.35 & 56,634 & 81.04 & 66.42 \\
		\textbf{\tool-Flat}   & \textbf{230} & \textbf{2.61} & \textbf{2.61} &  \textbf{25,238} & \textbf{33.81} & \textbf{47.75} \\
		\textbf{\tool}   & \textbf{231} & \textbf{0.43} & \textbf{0.00} &\textbf{33,340} & \textbf{4.09} & \textbf{0.02} \\
		%\hline
		\bottomrule
	\end{tabular}}
\end{table}

From the perspective of incompatibility issues, it can be seen from \Cref{tab:incom} that \tool reported 33,340 incompatibility issues with 4.09\% FP rate and 0.02\% FN rate, which demonstrated the ability of \tool to detect license incompatibility issues. We can also observe that the FP and FN rates of four baselines, ranging from 36.22\% to 99.28\%, are much higher than those of \tool, ranging from 0.02\% to 4.09\%. By analyzing the results, we found that the false negatives of SPDX-VT, and Librariesio were mainly due to their limited scope of licenses that can be handled. LiDetector localized more incompatibility issues than the other three baselines, due to the flexibility adaption to arbitrary licenses. However, it treated all component licenses equally, and ignored the potential hierarchy among them, which contributed to the false negatives. 
%SPDX-VT is similar as LiDetector and any pair of licenses with incompatibility relation input into its graph can lead to a incompatibility issue. 



\noindent $\bullet$ \textbf{Effects of hierarchy extraction}. 
To illustrate the usefulness of the license hierarchy on incompatibility detection, we also implemented \tool without hierarchy extraction, denoted by \tool-Flat in \Cref{tab:incom}. It can be seen that without hierarchy extraction, \tool reported 230 incompatible projects with 6 false positives and 6 false negatives,
% However, it has the same FN rate with LiDetector (i.e., 4.62\%). 
similar to the results of LiDetector.
The reason behind is that both \tool-Flat and LiDetector treated all component licenses in modules or accompanied with third-party packages equally important regardless of their licensing scopes. To be specific, \tool-Flat reported a project with licence incompatibility issues when there existed a component license that was more {restrictive} than the project license. However, \tool reported an incompatible project when there existed a child license more restrictive than its parent license. It can be inferred that given the same project and the same results of license understanding, the true positives found by \tool-Flat can theoretically be found by \tool. Nevertheless, when the parent license is not the project license, but a license in a module of the project, the parent and child licenses are regarded as two independent licenses by \tool-Flat and LiDetector, which might lead to a false negative. On the other hand, it can also be seen that the false positive rate of the issues reported by \tool-Flat is also higher than that reported by \tool. By observing the results, we found that it was due to the projects without project licenses. In this case, \tool-Flat considered all rights reserved, and created the most strict project license for incompatibility detection, which resulted in the false positives. %Moreover, it can also be seen that \tool-Flat reported fewer issues than \tool by {24.30\%}, since it ignored the incompatibility issues occurred between a parent-child pair of licenses where the parent license is not the project license.

\iffalse
\begin{table}[htp]
	\centering
	\footnotesize
	\caption{Comparison on License Incompatibility Localization} \label{tab:incom_comp}
	%\scalebox{0.9}{
	\begin{tabular}{ccccc}
		\toprule
		%\hline
		%\textbf{Tool} & \textbf{\#Pro.} & \textbf{\#Incomp. Pro.} & \textbf{\#Incomp. Issue} & \textbf{FP} & \textbf{FN}\\
		\textbf{Tool} & \makecell[c]{\textbf{\#Incomp. Pro.} \\\textbf{Reported}} & \makecell[c]{\textbf{\#Incomp.} \\ \textbf{Issues}} & \textbf{FP} & \textbf{FN}\\
		\hline
		Ninka   & 136 (136 overlap) & 32,289 (4,339 overlap) & 0 (0.00\%) & 102 (42.86\%)  \\
		SPDX-VT   & 76 (76 overlap) & - & 0 (0.00\%) & 162 (68.07\%)  \\
		Librariesio   & 171 (157 overlap) & - & 14 (8.19\%) & 81 (34.03\%)  \\
		LiDetector   & 228 (227 overlap) & 56,634 (9,835 overlap) & 1 (0.44\%) & 11 (4.62\%)  \\
		\textbf{\tool (flatten)}   & \textbf{227} & \textbf{90,556 (17,829 overlap)} & \textbf{0} (\textbf{0.00\%}) & \textbf{11} (\textbf{4.62\%})  \\
		\textbf{\tool (hierarchy)}   & \textbf{238} & \textbf{29,823} & \textbf{0} (\textbf{0.00\%}) & \textbf{0} (\textbf{0.00\%}) \\
		%\hline
		\bottomrule
	\end{tabular}%}
\end{table}
\fi


\noindent \fbox{
\parbox{0.95\linewidth}{
\textbf{Answer to RQ2:} 
\tool can effectively identify OSS with license incompatibility issues, achieving 0.43\% FP rate and no false negatives for 230 real-world incompatible OSS, outperforming four state-of-the-art approaches.
}}


\subsection{Evaluation on Incompatibility Resolution}
%In this section, we evaluate the performance of \tool on resolving license incompatibility issues. 
\subsubsection{{Setup}}
The experiment on license incompatibility resolution was conducted based on the same dataset with \Cref{sec:incomp}, which consist of {230} incompatible projects detected and verified at the previous step. For each project, there exists at least one pair of licenses which are incompatible with each other. Note that a project was regarded as resolved if all incompatibility issues could be resolved by \tool, and no new incompatibility issues were introduced into the resolved project according to \tool. In addition, to verify the correctness of resolutions provided by \tool, we also randomly selected {5\%} of the licenses recommended/generated by \tool (i.e., {1,484} licenses) for evaluation. %Note that as mentioned in \Cref{sec:custom}, we only generate a custom license when no existing official licenses can satisfy the constraints. Similarly, exceptions are attached only when the conflict licenses can not be modified by the software owner. %We conduct the evaluation which first checked whether exists an official license that can be used to resolve the current issue. Only when none of official licenses satisfied the license constraints, \tool generated a custom one as an alternative solution. Finally, {only when there were incompatible conflicts between multiple child licenses of the same parent license, \tool generated a custom license with an exception.} Note that \tool checked the copyright holders of incompatible child licenses. If a child license can be modified or changed to remove the conflict, \tool did not generate the exception.


\subsubsection{{Results}}

\begin{table}[]
    \centering\footnotesize
    \caption{License Incompatibility Resolved by \tool}
    \begin{tabular}{lrr}
    \toprule
        {\#Incompatible projects} &  230 & 100.00\%\\ 
        {\space \space-- \#Resolved projects}  &  \textbf{144} & \textbf{62.61\%}\\
        % {\space \space-- \#Unresolved incompatible projects}  &  {86} ({37.39\%})\\
        {\space \space-- \#Projects that cannot be resolved}  &  {86} & {37.39\%}\\
        % {\space \space-- \#Wrong resolutions provided by \tool}  &  {2} ({0.87\%})\\
        \midrule
        {\#Incompatible issues} &  33,340 & 100.00\%\\
        % {\#Incompatible issues found by \tool} &  \sihan{33,340} (100.00\%)\\
        {\space \space-- \#Issues resolved by recommending official licenses} &  \textbf{11,509} & \textbf{34.52\%}\\
        {\space \space-- \#Issues resolved by generating custom licenses} &  \textbf{18,185} & \textbf{54.54\%}\\
        {\space \space\space\space\space\space-- \#Custom licenses generated w/o exceptions} & 13,858  & 41.56\%\\
        {\space \space\space\space\space\space-- \#Custom licenses generated w/ exceptions} &  4,327 & 12.98\%\\
        %{$\bullet$ \#Issues resolved by custom licenses w/ exceptions} &  XXX (XXX\%)\\
        {\space \space-- \#Issues that cannot be resolved} &  3,646 & 10.94\%\\
        %{\space \space-- \#Wrong resolutions provided by \tool} &  514 (1.61\%)\\
        %{$\bullet$ \#Incompatible issues missed by \tool} &  \sihan{XXX} (XXX\%)\\
        \bottomrule
    \end{tabular}
    \label{tab:resolve}
\end{table}

\Cref{tab:resolve} shows the experimental results of \tool on license incompatibility resolution. Among the 230 incompatible projects, \tool resolved 144 projects, accounting for 62.61\% of all incompatible projects. From \Cref{tab:resolve} it can also be seen that {11,509} out of 33,340 incompatibility issues were resolved by \tool through recommending official licenses, accounting for {34.52\%} of all incompatibility issues. %Actually, there might exist more than one official license that satisfied the constraints. In this case, \tool produced a list of official licenses that can be used to replace the original license, and ranked according to their similarities to the original license via cosine similarity. 
%\sihan{Note that \tool provided resolutions to the problematic licenses one by one, since the choice of the solutions to one issue might affect the incompatibility detection and resolution for other licenses.}
%
%It can also be seen that official licenses could only resolve a small portion of license incompatibility issues. 
Since a majority of the incompatibility issues cannot be resolved via official licenses, as an alternative way, \tool also resolved issues by generating custom licenses to satisfy the constraints. In total, \tool generated {18,185} custom licenses, among which only {4,327} custom licenses were attached with specific exceptions to address with the problem when multiple child licenses of a target license are conflict with each other, and none of them can be modified by the software owner to resolve the conflict. 


From \Cref{tab:resolve} it can also be observed that there are also some license incompatibility issues that cannot be resolved by \tool. The reasons are two folds. First, as mentioned in~\Cref{sec:constraint}, not all constraints from the context of a given license can be resolved; there exists some cases where the parent license of the current license is less restrictive than a child license of the current license. Second, not all licenses are allowed to be modified or replaced by the software owner. Some licenses were incorporated along with third-party software packages, and the project owner has no rights to change them. %To filter out licenses that can not be modified by the software owner, \tool detected copyright holders for all licenses involved in incompatibility issues and determined the licenses that can be modified. 
In total, {3,646} incompatible issues could not be resolved by \tool, accounting for {10.94\%} of the license incompatibility issues detected by \tool at the previous step.

Finally, to verify the correctness of resolutions provided by \tool, we randomly selected 5\% of the {29,694} licenses recommended or generated by \tool. We manually verified and cross-validated the correctness of these licenses. 
% We found that provided {510} wrong resolutions for the incompatible projects, accounting for {1.53\%} of all issues. 
By analyzing the results, {we found that {1.48\%} of verified resolutions were incorrect, mainly caused by the incorrect predictions in the phases of fine-grained entity extraction and relation extraction, which affected the license constraints for recommending official licenses and the detailed information used for generating custom licenses. }



\noindent $\bullet$ \textbf{Effects of copyright holder}. \Cref{tab:copyright} displays the results of license incompatibility resolution when allowing different ranges of licenses to be modified. In this paper, we classify licenses in a project into the project license and other licenses (also known as component licenses). Unlike previous studies that only focus on the project license, this paper observed a considerable number of component licenses which were also written by the project owner and thus could be modified for resolving license incompatibility. To investigate how the ranges of the licenses allowed to be modified affects the results of license incompatibility resolution, we implemented \tool with different settings. The experiment was conducted on the same dataset with \Cref{tab:resolve}.

From \Cref{tab:copyright} it can be seen that %when all licenses were allowed to be changed,
% \tool resolved \sihan{231} problematic projects and \sihan{33,340} incompatible issues, achieving \sihan{100\%} and \sihan{100\%} of \gy{the accuracy}, respectively. 
%\tool resolved 100\% of problematic projects and 100\% of incompatible issues. 
% By observing the results, we found that \sihan{XX\%} of the incompatibility issues were resolved by modifying the component licenses. 
%However, since many licenses are introduced accompanied with third-party software packages or modules, which could not be modified by the project owner, we also evaluated the performance of \tool when only project licenses in a given project were allowed to be changed. We can see that 
{138} projects and {26,523} incompatibility issues could be handled with only the project licenses allowed to be modified, accounting for {60.00\%} and {79.55\%} of the incompatible projects and issues, respectively. %The results indicate that more than half of the license incompatibility issues can be resolved by the project licenses. %the limitation to project licenses might contribute to many unsolved projects and issues. 
However, when only allowing the licenses with the same copyright holder with the project license to be changed, \tool only resolved {76} projects and {23,859} issues. By analyzing the results, we found that many project licenses did not explicitly state their copyright holders in their license texts, resulting in the low success rate of resolving incompatibility issues. 
%
In this paper, we propose to modify two types of licenses (i.e., the project licenses and the component licenses with the same copyright holder). The first intuition is that the project owner is ought to have the rights to modify or replace the project license, whose licensing scope is the whole project. The second intuition is that only a component license states the same copyright holder with as the project license states, the project owner should be allowed to modify it. From \Cref{tab:copyright} it can be seen that \tool resolved 2.61\% more projects and 9.51\% more incompatibility issues by adding copyright holder detection.
%The results also show that \tool resolved \sihan{XX} more issues by modifying component licenses with the same copyright holder with the project license.

\begin{table}[t]
	\centering
	\footnotesize
	\caption{Effects of Copyright Holder Detection} \label{tab:copyright}
	%\scalebox{0.9}{
	\begin{tabular}{ccc|cc}
		\toprule
		%\hline
		%\textbf{Tool} & \textbf{\#Pro.} & \textbf{\#Incomp. Pro.} & \textbf{\#Incomp. Issue} & \textbf{FP} & \textbf{FN}\\
		\multirow{2}{*}{\textbf{Modified}} & \multicolumn{2}{c|}{\textbf{\#Incomp. Projects}} & \multicolumn{2}{c}{\textbf{\#Incomp. Issues}}\\\cline{2-5}
		 & {\textbf{\#Resolved Pro.}} & \textbf{Suc. Rate} & \textbf{\#Resolved Issues}& \textbf{Suc. Rate} \\
		\hline
% 		\textbf{ALL}   & 231 & 100.00 & 33,340 & 100.00  \\
		{PL}   & 138 & 60.00\% & 26,523 & 79.55\%  \\
		{CH}   & 76 & 33.04\% & 23,859 & 71.56\%  \\
	  \textbf{PL+CH}  & \textbf{144} & \textbf{62.61\%} & \textbf{29,694} & \textbf{89.06\%}  \\
		%\hline
		\bottomrule
	\end{tabular}%}
%\vspace{-1.5mm}
\begin{center}
	    \footnotesize
	    %\textit{\textbf{ALL}: All licenses involved in license incompatibly issues can changed.} \\
	    {\textit{\textbf{PL}: Only the project license can be changed to resolve license incompatibility.}\\
	    \textit{\textbf{CH}: Only licenses whose copyright holder is the project owner can be changed.}}
\end{center}
\end{table}

\noindent \fbox{
\parbox{0.95\linewidth}{
\textbf{Answer to RQ3:} 
By recommending official licenses in priority and generating custom licenses as alternative solutions, \tool resolved 144 incompatible projects, among which 34.52\% were resolved by official licenses, 41.56\% and 12.98\% were resolved by custom licenses w/o and w/ exceptions, respectively.
}}




\begin{comment}
\begin{table}[]
    \centering\small
    \caption{Caption}
    \label{tab:my_label1}
    \begin{tabular}{cc}
        \toprule
        \textbf{Incompatible projects} & \textbf{Successfully resolved projects} \\
        \midrule
        & \\
        \bottomrule
    \end{tabular}
\end{table}

\begin{table}[]
    \centering\small
    \caption{Caption}
    \label{tab:my_label2}
    \begin{tabular}{cccc}
        \toprule
        \textbf{\#Incomp. issues} & \textbf{\#Resol. by O} & \textbf{Resol. by C} & \textbf{Unresol.}\\
        \midrule
         & & & \\
        \bottomrule
    \end{tabular}
\end{table}


\begin{table}[]
\small
\caption{xxx}
\begin{tabular}{c|p{2cm}<{\centering}c}
\toprule
\textbf{\#Incompatible projects}                & \multicolumn{2}{c}{\textbf{\#Successfully resolved projects}} \\ \hline
XXX                                             & \multicolumn{2}{c}{XXX}                                       \\ \midrule
\multirow{2}{*}{\textbf{\#Incompatible issues}} & \multicolumn{2}{c}{\textbf{\#Successfully resolved issues}}   \\ \cline{2-3} 
    & \textbf{Official}    & \textbf{Custom}    \\ \hline
XXX   & XXX     & XXX                 \\ 
\bottomrule
\end{tabular}
\end{table}
\end{comment}






\iffalse

\begin{table}
	\centering
	\small
	\caption{Comparison on Entity Extraction}
	\scalebox{0.9}{\begin{tabular}{ccccc}
		\toprule
		%\hline
		\textbf{Method} & \textbf{Act (\%)} & \textbf{Rec (\%)} & \textbf{Att (\%)}& \textbf{Con (\%)}\\
		\midrule
		Regular Matching  & XX/XX  & XX/XX & XX/XX  & XX/XX\\
		Semantic Similarity  & XX/XX  & XX/XX & XX/XX  & XX/XX\\
		FOSS-LTE  & XX/XX  & -/- & XX/XX  & -/-\\
		LiDetector & XX/XX  & -/- & XX/XX  & XX/XX\\
		\tool  & XX/XX  & XX/XX & XX/XX  & XX/XX\\
		%\hline
		\bottomrule
		\label{tab:entity}
	\end{tabular}}
\end{table}
\fi